\documentclass[11pt]{article}
\usepackage[T1]{fontenc}
\usepackage[utf8]{inputenc}
\usepackage{lmodern}
\usepackage{amsmath,amssymb,amsthm,mathtools}
\usepackage{geometry}
\usepackage{microtype}
\usepackage{booktabs,longtable,array}
\usepackage{enumitem}
\usepackage{xcolor}
\usepackage{listings}
\usepackage{hyperref}
\usepackage{bookmark}
\usepackage{url}
\geometry{margin=1in}
\hypersetup{colorlinks=true,linkcolor=blue!45!black,citecolor=green!35!black,urlcolor=blue!55!black,pdftitle={Alaoglu--Erdos Conjecture: Continued Investigation},pdfauthor={OpenAI GPT-5.6 Pro}}
\allowdisplaybreaks
\setlength{\parindent}{0pt}
\setlength{\parskip}{0.55em}
\setlist[itemize]{topsep=0.25em,itemsep=0.2em}
\setlist[enumerate]{topsep=0.25em,itemsep=0.2em}

\newtheorem{theorem}{Theorem}[section]
\newtheorem{lemma}[theorem]{Lemma}
\newtheorem{corollary}[theorem]{Corollary}
\newtheorem{proposition}[theorem]{Proposition}
\newtheorem{conjecture}[theorem]{Conjecture}
\theoremstyle{definition}
\newtheorem{definition}[theorem]{Definition}
\newtheorem{remark}[theorem]{Remark}
\newtheorem{experiment}[theorem]{Computational experiment}

\newcommand{\N}{\mathbb N}
\newcommand{\Nzero}{\mathbb N_0}
\newcommand{\Z}{\mathbb Z}
\newcommand{\Q}{\mathbb Q}
\newcommand{\R}{\mathbb R}
\newcommand{\C}{\mathbb C}
\newcommand{\vp}{v_p}
\newcommand{\ee}{\mathrm e}
\newcommand{\ord}{\operatorname{ord}}
\newcommand{\num}{\operatorname{num}}
\newcommand{\den}{\operatorname{den}}
\newcommand{\rad}{\operatorname{rad}}
\newcommand{\Span}{\operatorname{span}}
\newcommand{\Law}{\operatorname{Law}}
\newcommand{\supp}{\operatorname{supp}}
\newcommand{\vol}{\operatorname{vol}}
\newcommand{\cO}{\mathcal O}
\newcommand{\DeltaV}{\mathop{\Delta}\nolimits}
\newcommand{\E}{\mathbb E}
\newcommand{\Var}{\operatorname{Var}}
\newcommand{\Cov}{\operatorname{Cov}}
\newcommand{\dist}{\operatorname{dist}}
\newcommand{\AE}{Alaoglu--Erd\H{o}s}

\title{Toward the Alaoglu--Erd\H{o}s Exponent Conjecture\\
\large Exact solution-monoid rigidity, adelic interpolation determinants,\\
and a new triangular-grid exclusion inequality}
\author{Continuation and audit of the attached unified report}
\date{August 2026}


\lstdefinestyle{proofpython}{
  language=Python,
  basicstyle=\ttfamily\scriptsize,
  columns=fullflexible,
  keepspaces=true,
  breaklines=true,
  breakatwhitespace=false,
  frame=single,
  showstringspaces=false,
  tabsize=4
}

\begin{document}
\maketitle

\begin{abstract}
We continue the investigation of the conjecture that a real number $x$ satisfying
$2^x,3^x\in\Z$ must itself be an integer.  The attached unified report had
already reduced the problem to a very rigid hypothetical counterexample and
had identified interpolation determinants as a promising, but critically
balanced, route.  This continuation first audits the structural reduction in
full detail.  Using the Six Exponentials Theorem, local finiteness, and a
valuation argument, we prove that if any counterexample exists then there is a
unique least one, denoted $\beta$, and the complete solution set is exactly the
free additive monoid
\[
   \{x\ge 0:2^x,3^x\in\Z\}=\Nzero+\Nzero\beta.
\]
Thus one counterexample would not produce a miscellaneous exceptional set: it
would force an exact two-dimensional lattice of integral values.

We then retain that two-dimensional structure in a family of generalized
Vandermonde determinants.  Their archimedean size is controlled by total
positivity and the Harish--Chandra--Itzykson--Zuber integral, while their finite
prime divisibility is bounded below by optimal-assignment costs for the
valuation matrix.  For square grids this gives the explicit divisor
$6^{(1+\beta)L_m}$, where
$L_m=m^2(m-1)(m-2)/6$.  More efficiently, using the first $N$ lattice nodes in
triangular order gives a rigorous necessary inequality
\[
  g\!\left(\sqrt{\beta\log 2\log 3}\right)\ge 0,
  \qquad
  g(z)=\frac12\log(2z)-1+\left(\frac{11}{3}-\pi\right)z.
\]
Its unique positive zero is
$z_0=1.12895010749713$, corresponding to
$\beta_\triangle=1.67370758736671$.  In particular the hypothetical least
counterexample cannot have $2^\beta=3$; combined with integrality and
nonintegrality this yields the unconditional reduction $2^\beta\ge5$ and hence
$\beta\ge\log_2 5$.

Finally, we derive an adelic version valid after any unimodular rational change
of the bases $2,3$.  It makes the denominator cost completely explicit and
turns continued-fraction changes of basis into rigorous gcd and denominator
constraints on $A=2^\beta$ and $B=3^\beta$.  The resulting framework does not
yet prove the full conjecture.  It does, however, identify precisely what a
successful determinant proof still needs: an $N^2$-scale finite-prime gain
beyond the componentwise transport bound, most plausibly from joint valuation
directions, confluent integer determinants, or a variational choice of lattice
regions.


A proof-producing finite verification, using directed interval arithmetic and
an explicit logarithmic-series remainder, additionally rules out every
$A=2^\beta\le100{,}000$ except the known power-of-two cases.  Thus any
counterexample would have $A>100{,}000$ and
$\beta>\log_2(100{,}000)\approx16.6096404744$.


No complete unconditional proof of the conjecture is claimed here.  Every new
statement is labelled according to whether it is proved, computationally
certified, conditional, or proposed as a research direction.
\end{abstract}

\tableofcontents
\clearpage

\section{Scope, status, and relation to the attached report}

The problem is
\begin{conjecture}[\AE\ exponent conjecture]\label{conj:AE}
If $x\in\R$ and $2^x,3^x\in\Z$, then $x\in\Z$.
\end{conjecture}

This document is a sequel rather than a replacement for the attached unified
report.  It restates every inherited fact used below, because several later
arguments depend sensitively on the exact quantifiers and on whether a
quantity is integral or merely rational.  In particular, the slogan ``one
counterexample generates a grid'' is proved here as an exact equality of
monoids, not used as a heuristic.

No complete proof of Conjecture~\ref{conj:AE} is claimed.  The strongest new
unconditional outputs of this continuation are:
\begin{enumerate}
\item a streamlined proof of the exact free-monoid structure, including the
      denominator-elimination step that is easy to understate;
\item a universal $N^2$-scale prime-divisibility lower bound for exact-grid
      interpolation determinants;
\item the triangular determinant inequality of
      Theorem~\ref{thm:triangle-inequality};
\item a rational-basis, product-formula extension that quantifies exactly how
      denominators defeat naive lattice reduction;
\item additional shared-prime, congruence, counting, and formalization
      consequences.
\end{enumerate}

\subsection{Claim ledger}

\begin{center}
\begin{tabular}{p{0.43\textwidth}p{0.16\textwidth}p{0.31\textwidth}}
\toprule
Claim & Status & Main input \\
\midrule
A counterexample is positive and transcendental & proved & elementary reductions; Gelfond--Schneider \\
All solutions lie in a rational affine line through $1$ and one nonintegral solution & proved & Six Exponentials Theorem \\
The complete solution set is $\Nzero+\Nzero\beta$ & proved & valuations and minimality \\
A counterexample has finite effective irrationality measure & proved & linear forms in two logarithms \\
Universal square-grid divisor $6^{(1+\beta)L_m}$ & proved & tropical determinant valuation \\
Triangular inequality $g(\sqrt{\beta\log2\log3})\ge0$ & proved & HCIZ bound, assignment asymptotics, Baker uniformity \\
Rational-basis denominator inequality & proved & product formula and signed assignment bounds \\
A complete contradiction for all $\beta$ & not obtained & requires further $N^2$-scale arithmetic gain \\
\bottomrule
\end{tabular}
\end{center}

\section{Elementary reductions and notation}

Set
\[
  \mathcal S:=\{x\in\R:2^x\in\Z,\ 3^x\in\Z\}.
\]
Since $2^x,3^x>0$, the integers in the definition are positive.  Write
\[
  \alpha:=\log 2,\qquad \gamma:=\log 3.
\]
All logarithms in the report are natural unless a base is displayed.

\begin{lemma}[No negative or small positive solutions]\label{lem:range}
One has $\mathcal S\subseteq[0,\infty)$, and
$\mathcal S\cap(0,1)=\varnothing$.
\end{lemma}
\begin{proof}
If $x<0$, then $0<2^x<1$, so $2^x$ is not an integer.  If $0<x<1$, then
$1<2^x<2$, again impossible.
\end{proof}

\begin{lemma}[Rational solutions are integral]\label{lem:rational}
If $x\in\Q$ and $2^x\in\Z$, then $x\in\Z$.  Consequently every rational member
of $\mathcal S$ is a nonnegative integer.
\end{lemma}
\begin{proof}
Write $x=a/d$ with $d>0$ and $\gcd(a,d)=1$.  If $2^{a/d}=M\in\Z$, then
$M^d=2^a$.  Unique factorization gives $d\mid a$, hence $d=1$.  The range
statement follows from Lemma~\ref{lem:range}.
\end{proof}

\begin{corollary}[A counterexample is transcendental]\label{cor:transcendental}
Every nonintegral member of $\mathcal S$ is transcendental.
\end{corollary}
\begin{proof}
It is irrational by Lemma~\ref{lem:rational}.  If it were algebraic irrational,
Gelfond--Schneider would make $2^x$ transcendental, contradicting
$2^x\in\Z$.
\end{proof}

\begin{lemma}[Local finiteness]\label{lem:local-finite}
For every $T\ge0$, the set $\mathcal S\cap[0,T]$ is finite.
\end{lemma}
\begin{proof}
The map $x\mapsto2^x$ is injective, and its values on this set are integers in
$[1,2^T]$.  Hence there are at most $\lfloor2^T\rfloor$ possibilities.
\end{proof}

The set $\mathcal S$ is an additive submonoid of $\R_{\ge0}$: if $x,y\in
\mathcal S$, then $2^{x+y}=2^x2^y$ and $3^{x+y}=3^x3^y$ are integers.  It
contains $\Nzero$.

\section{The exact one-counterexample structure}

This section audits and sharpens the strongest structural conclusion of the
attached report.

\subsection{Affine rank at most two}

We use the following standard transcendence theorem.

\begin{theorem}[Six Exponentials Theorem]\label{thm:six-exp}
Let $u_1,u_2\in\C$ be linearly independent over $\Q$, and let
$v_1,v_2,v_3\in\C$ be linearly independent over $\Q$.  Then at least one of
$\exp(u_i v_j)$, $1\le i\le2$, $1\le j\le3$, is transcendental.
\end{theorem}

\begin{proposition}[Affine rank theorem]\label{prop:affine-rank}
If $x,y\in\mathcal S$, then $1,x,y$ are linearly dependent over $\Q$.
Consequently, if $\beta\in\mathcal S\setminus\Q$ is fixed, every
$x\in\mathcal S$ has a unique representation
\[
  x=r+s\beta,\qquad r,s\in\Q.
\]
\end{proposition}
\begin{proof}
The numbers $\alpha=\log2$ and $\gamma=\log3$ are $\Q$-linearly independent,
since $2^a3^b=1$ with $a,b\in\Z$ forces $a=b=0$.  If $1,x,y$ were
$\Q$-linearly independent, Theorem~\ref{thm:six-exp} applied to
$(\alpha,\gamma)$ and $(1,x,y)$ would say that at least one of
\[
  2,\ 3,\ 2^x,\ 3^x,\ 2^y,\ 3^y
\]
is transcendental.  All six are integers, a contradiction.  For the second
claim, the dependence of $1,\beta,x$ has a nonzero coefficient of $x$, because
$1,\beta$ are independent.  Uniqueness follows from the same independence.
\end{proof}

Thus Six Exponentials already rules out two genuinely independent exceptional
directions.  What remains is to eliminate rational denominators in the affine
coordinates.  That step uses minimality and prime valuations.

\subsection{The primitive least counterexample}

Assume from now until explicitly stated otherwise that Conjecture~\ref{conj:AE}
is false.  By Lemma~\ref{lem:local-finite}, there is a least nonintegral member
\[
  \beta:=\min(\mathcal S\setminus\Nzero).
\]
By Lemma~\ref{lem:range}, $\beta>1$.  Put
\[
  A:=2^\beta\in\N,\qquad B:=3^\beta\in\N.
\]
The number $\beta$ is irrational and transcendental, and $A$ is not a power of
$2$, while $B$ is not a power of $3$.

For each prime $p$, write
\[
  a_p:=v_p(A),\qquad b_p:=v_p(B).
\]
Only finitely many of these numbers are nonzero.

\begin{lemma}[Shifted perfect-power primitivity]\label{lem:primitive-gcd}
The integer
\[
 g_0:=\gcd\!\left(
   \{a_p:p\ne2\}\cup\{b_p:p\ne3\}\cup\{a_2-b_3\}
 \right)
\]
(where signs are ignored in taking the gcd) is equal to $1$.
\end{lemma}
\begin{proof}
Suppose a prime $q$ divides every displayed integer.  Let $r$ be the common
residue of $a_2$ and $b_3$ modulo $q$, chosen with $0\le r<q$.  Then
$a_2-r$ and $b_3-r$ are nonnegative multiples of $q$, and all other relevant
valuations are multiples of $q$.  Hence there are integers $U,V>0$ such that
\[
  A=2^rU^q,\qquad B=3^rV^q.
\]
Therefore
\[
  2^{(\beta-r)/q}=U,\qquad 3^{(\beta-r)/q}=V.
\]
The exponent $(\beta-r)/q$ cannot be negative, because its power of $2$ is a
positive integer; it cannot be zero, because $\beta$ is irrational.  It is thus
a positive nonintegral member of $\mathcal S$, and it is strictly smaller than
$\beta$, contradicting minimality.
\end{proof}

\begin{lemma}[Denominator elimination]\label{lem:denominator-elimination}
If $x\in\mathcal S$ and
\[
  x=\frac{a+b\beta}{d},\qquad a,b,d\in\Z,\quad d>0,
\]
with $\gcd(a,b,d)=1$, then $d=1$.
\end{lemma}
\begin{proof}
Let $C=2^x$ and $D=3^x$.  Then
\[
  C^d=2^aA^b,\qquad D^d=3^aB^b.
\]
Taking $p$-adic valuations, including when $a$ or $b$ is negative, gives
\begin{align*}
  d\,v_p(C)&=a\,\mathbf 1_{p=2}+b a_p,\\
  d\,v_p(D)&=a\,\mathbf 1_{p=3}+b b_p.
\end{align*}
Consequently $d$ divides $b a_p$ for every $p\ne2$, divides $b b_p$ for every
$p\ne3$, and, by comparing the congruences at $p=2$ and $p=3$, divides
$b(a_2-b_3)$.  Lemma~\ref{lem:primitive-gcd} and B\'ezout therefore imply
$d\mid b$.  The congruence at $p=2$ then gives $d\mid a$.  Since
$\gcd(a,b,d)=1$, one has $d=1$.
\end{proof}

\subsection{Exact free-monoid theorem}

\begin{theorem}[Exact solution monoid]\label{thm:exact-monoid}
If Conjecture~\ref{conj:AE} is false and $\beta$ is the least nonintegral
solution defined above, then
\[
  \boxed{\mathcal S=\{n+k\beta:n,k\in\Nzero\}.}
\]
The representation is unique.  Moreover,
\[
  \min\{v_2(A),v_3(B)\}=0.
\]
\end{theorem}
\begin{proof}
Proposition~\ref{prop:affine-rank} and
Lemma~\ref{lem:denominator-elimination} show that every $x\in\mathcal S$ has a
unique representation $x=n+k\beta$ with $n,k\in\Z$.

Suppose $k<0$, and put $h=-k>0$.  The integrality of
$2^x=2^n/A^h$ forces $A$ to have no prime divisor other than $2$, while the
integrality of $3^x=3^n/B^h$ forces $B$ to have no prime divisor other than
$3$.  Then $A=2^r$ and $\beta=r\in\Z$, a contradiction.  Thus $k\ge0$.
If $k=0$, integrality and Lemma~\ref{lem:range} give $n\ge0$.

Let $k>0$.  All valuations away from $2$ and $3$ are automatically
nonnegative.  The two remaining conditions are
\[
  n+k v_2(A)\ge0,\qquad n+k v_3(B)\ge0.
\]
Put $c=\min\{v_2(A),v_3(B)\}$.  If $c>0$, then
$2^{\beta-c}=A/2^c$ and $3^{\beta-c}=B/3^c$ are integers.  As in the proof of
Lemma~\ref{lem:primitive-gcd}, the exponent $\beta-c$ is positive,
nonintegral, and smaller than $\beta$, impossible.  Hence $c=0$, and the two
inequalities force $n\ge0$.

Conversely every $n,k\ge0$ gives integer powers
$2^{n+k\beta}=2^nA^k$ and $3^{n+k\beta}=3^nB^k$.  Uniqueness follows from the
irrationality of $\beta$.
\end{proof}

\begin{corollary}[One base-prime unit]\label{cor:unit}
A least counterexample satisfies at least one of
\[
  2\nmid A,\qquad 3\nmid B.
\]
It also satisfies the shifted primitivity condition of
Lemma~\ref{lem:primitive-gcd}.
\end{corollary}

\begin{corollary}[Power-map monoid formulation]\label{cor:power-monoid}
Let $\rho=\log_2 3$.  If a counterexample exists, then for the least one
$A=2^\beta$, $B=3^\beta$,
\[
 \{m\in\N:m^\rho\in\N\}=\{2^nA^k:n,k\in\Nzero\},
\]
and the power map is the monoid isomorphism
\[
  2^nA^k\longmapsto3^nB^k.
\]
\end{corollary}
\begin{proof}
For $m\in\N$, put $x=\log_2m$.  Then $m^\rho=3^x$.  Thus
$m^\rho\in\N$ exactly when $x\in\mathcal S$, and
Theorem~\ref{thm:exact-monoid} applies.
\end{proof}

This formulation is useful conceptually: a counterexample would force the
integer points of the transcendental power map $m\mapsto m^{\log_2 3}$ to form
an exact free rank-two multiplicative monoid, rather than an accidental sparse
set.

\section{Effective Diophantine and counting consequences}

The two logarithmic descriptions of $\beta$ imply more than transcendence.
Because $A$ is not a power of $2$, the linear form
$q\log A-p\log2$ is nonzero whenever $(p,q)\ne(0,0)$.  Effective lower bounds
for linear forms in logarithms therefore give constants $c>0$ and $\tau>0$,
effectively computable from $A$, such that
\begin{equation}\label{eq:finite-type}
  \|q\beta\|\ge c q^{-\tau}\qquad(q\ge1).
\end{equation}
The analogous representation $\beta=\log B/\log3$ gives another such bound.
No useful small absolute exponent is asserted here; standard explicit bounds
are enormous.  The qualitative conclusion is nevertheless important.

\begin{corollary}[Finite irrationality measure]\label{cor:finite-mu}
A counterexample $\beta$ is not a Liouville number.  Its irrationality measure
is finite and effectively bounded in terms of either $A$ or $B$.
\end{corollary}

The exact monoid also gives an exact counting formula.  Let
\[
  N_{\mathcal S}(T):=\#(\mathcal S\cap[0,T]).
\]
Then, with $K=\lfloor T/\beta\rfloor$,
\begin{equation}\label{eq:count-exact}
 N_{\mathcal S}(T)=\sum_{k=0}^{K}\bigl(\lfloor T-k\beta\rfloor+1\bigr).
\end{equation}

\begin{proposition}[Quadratic counting with a power saving]\label{prop:count}
There is an effective $\delta>0$ such that
\[
 N_{\mathcal S}(T)=\frac{T^2}{2\beta}
 +\frac{\beta+1}{2\beta}T+O_\beta(T^{1-\delta}).
\]
\end{proposition}
\begin{proof}[Proof sketch]
Equation~\eqref{eq:finite-type}, the Erd\H{o}s--Tur\'an inequality, and the
geometric-series estimate for $\sum_{k\le K}e^{2\pi i h k\beta}$ give a power
saving for the discrepancy of the rotation sequence $\{k\beta\}$.  In
particular,
\[
 \sum_{k=0}^{K}\{T-k\beta\}=\frac{K+1}{2}+O(K^{1-\delta})
\]
for some effective $\delta>0$.  Insert this in
Equation~\eqref{eq:count-exact}.  Writing $K=T/\beta-\{T/\beta\}$, the
apparently oscillatory terms of order $T$ cancel, leaving the displayed linear
coefficient.
\end{proof}

The same argument gives quantitative equidistribution of the dense orbit
\[
 \left(2^{\{k\beta\}},3^{\{k\beta\}}\right)
 =\left(\frac{A^k}{2^{\lfloor k\beta\rfloor}},
         \frac{B^k}{3^{\lfloor k\beta\rfloor}}\right).
\]
Thus a counterexample would generate a dense family of rational points on the
arc $t\mapsto(2^t,3^t)$, with pure base-power denominators.  Their heights grow
exponentially in $k$, so this observation alone is far below the density needed
for Pila--Wilkie-type counting to force algebraicity; it is nevertheless a
useful reformulation for $S$-unit methods.


\section{Valuation-lattice saturation and an exact duality}
\label{sec:saturation-duality}

The denominator argument has a useful invariant formulation.  It also has a
symmetric consequence that is easy to miss if one works only with the least
counterexample $\beta$.

Write
\[
 a_p=v_p(A),\qquad b_p=v_p(B),
\]
and form the two-column integer matrix whose rows are
\[
  (\mathbf 1_{p=2},a_p)
  \quad\hbox{and}\quad
  (\mathbf 1_{p=3},b_p)
  \qquad (p\text{ prime}).
\]
Only finitely many entries in the second column are nonzero.  Let
$\mathcal L$ be the rank-two lattice generated by the columns in the direct
sum of all row copies of $\mathbb Z$.

\begin{theorem}[Saturated valuation lattice]
\label{thm:saturated-valuation-lattice}
The lattice $\mathcal L$ is primitive (saturated):
\[
  (\mathcal L\otimes_{\mathbb Z}\mathbb Q)
  \cap\bigoplus_p(\mathbb Z\oplus\mathbb Z)=\mathcal L.
\]
Equivalently, the greatest common divisor of all maximal minors of the matrix
is $1$.  In the present sparse matrix this is exactly
\[
 \gcd\!\left(
   \{a_p:p\ne2\}\cup
   \{b_p:p\ne3\}\cup
   \{a_2-b_3\}
 \right)=1.
\]
\end{theorem}

\begin{proof}
The nonzero maximal minors are, up to sign, exactly the displayed $a_p$,
$b_p$, and $a_2-b_3$: the first column is nonzero only in the row belonging to
$p=2$ in the first valuation block and the row belonging to $p=3$ in the
second.  It remains to prove that their gcd is $1$.

If a number $d>1$ divided all these minors, choose
$t\in\{0,\ldots,d-1\}$ with $t\equiv-a_2\pmod d$.  The congruence
$a_2\equiv b_3\pmod d$ then gives $t\equiv-b_3\pmod d$, while every other
valuation of $A$ or $B$ is already divisible by $d$.  Consequently both
$2^tA$ and $3^tB$ are perfect $d$th powers in $\mathbb N$.  Hence
\[
  \frac{\beta+t}{d}\in S.
\]
This number is nonintegral, and it is strictly between $0$ and $\beta$ because
$\beta>1$ and $0\le t\le d-1$, contradicting the choice of $\beta$.  Thus the
gcd of the maximal minors is $1$.  The rank-two Smith normal form criterion
then says exactly that the column lattice is saturated.
\end{proof}

A particularly concrete form of the theorem records exactly how extraction of
common roots behaves throughout the solution monoid.

\begin{corollary}[Perfect-power content is preserved]
\label{cor:perfect-power-content}
For $u,v\in\mathbb N_0$, not both zero, define
\[
 X_{u,v}=2^uA^v,\qquad Y_{u,v}=3^uB^v.
\]
Then
\[
 \gcd\!\left(
   \{v_p(X_{u,v}):p\text{ prime}\}
   \cup
   \{v_p(Y_{u,v}):p\text{ prime}\}
 \right)=\gcd(u,v).
\]
Equivalently, for every $d\ge1$, the two integers $X_{u,v}$ and $Y_{u,v}$
are simultaneously perfect $d$th powers if and only if
$d\mid u$ and $d\mid v$.
\end{corollary}

\begin{proof}
Every divisor of $\gcd(u,v)$ plainly divides all displayed valuations.  In the
other direction, if $d$ divides all of them, then
$X_{u,v}=C^d$ and $Y_{u,v}=D^d$ for positive integers $C,D$.  Thus
$(u+v\beta)/d\in S$.  Uniqueness of the representation in the irrational
basis $(1,\beta)$ gives $u/d,v/d\in\mathbb N_0$.
\end{proof}

For example, $A$ and $B$ cannot themselves be simultaneous proper powers; and
for every $t\ge0$, neither pair
\[
  (2^tA,3^tB),\qquad (2A^t,3B^t)
\]
can consist of simultaneous proper powers of the same degree.  These are not
independent ad hoc restrictions: all of them are specializations of one
saturated rank-two valuation lattice.

There is also an exact reciprocal description.  Put
\[
 S_{A,B}=\{y\in\mathbb R:A^y\in\mathbb Z,
                              \ B^y\in\mathbb Z\}.
\]

\begin{theorem}[Reciprocal solution monoid]
\label{thm:reciprocal-monoid}
Under the counterexample hypothesis,
\[
  S_{A,B}=\mathbb N_0+\frac1\beta\mathbb N_0.
\]
In particular $1/\beta$ is the least positive nonintegral element of
$S_{A,B}$.
\end{theorem}

\begin{proof}
If $y\in S_{A,B}$, then $x=\beta y$ satisfies
\[
  2^x=A^y\in\mathbb Z,
  \qquad
  3^x=B^y\in\mathbb Z.
\]
Hence $x=n+k\beta$ with $n,k\in\mathbb N_0$, and
$y=k+n/\beta$.  Conversely,
\[
 A^{k+n/\beta}=A^k2^n,
 \qquad
 B^{k+n/\beta}=B^k3^n
\]
are integers for all $n,k\ge0$.
\end{proof}

This reciprocal monoid shows that the two columns of the valuation matrix may
be interchanged without losing arithmetic information: the original and dual
root-extraction problems are the same saturated lattice viewed in opposite
bases.


\section{Exact-grid interpolation determinants}

The free monoid in Theorem~\ref{thm:exact-monoid} supplies two exact
rank-two grids: one in the logarithms of the bases and one in the solution
exponents.  Retaining both coordinates is the central step of this
continuation.

For finite sets $\mathcal R,\mathcal C\subset\Nzero^2$ with
$|\mathcal R|=|\mathcal C|=N$, define
\begin{align*}
 u_{a,b}&:=a\alpha+b\gamma,\qquad (a,b)\in\mathcal R,\\
 v_{n,k}&:=n+k\beta,\qquad (n,k)\in\mathcal C,
\end{align*}
and the matrix
\begin{equation}\label{eq:grid-matrix}
 M_{(a,b),(n,k)}
 :=\exp(u_{a,b}v_{n,k})
 =2^{an}3^{bn}A^{ak}B^{bk}.
\end{equation}
Every entry is a positive integer.

\subsection{Nonvanishing and total positivity}

\begin{lemma}[Strict positivity]\label{lem:tp}
Order the $u$-values and the $v$-values increasingly.  Then
\[
 D(\mathcal R,\mathcal C):=\det\bigl(\exp(u_i v_j)\bigr)_{i,j=1}^N>0.
\]
In particular it is a positive integer and hence at least $1$.
\end{lemma}
\begin{proof}
The $u$-values are distinct because $2$ and $3$ are multiplicatively
independent, and the $v$-values are distinct because $\beta$ is irrational.
The kernel $K(u,v)=e^{uv}$ is strictly totally positive on $\R\times\R$.
Equivalently, the functions $e^{v_1u},\ldots,e^{v_Nu}$ with
$v_1<\cdots<v_N$ form an extended complete Chebyshev system: their Wronskian is
\[
 e^{u(v_1+\cdots+v_N)}\prod_{i<j}(v_j-v_i)>0.
\]
The generalized mean-value theorem for Chebyshev systems gives the desired
sign.
\end{proof}

\subsection{Archimedean upper bound}

Write $\DeltaV(u)=\prod_{i<j}(u_j-u_i)$ and similarly for $v$, and let
$G(N+1)=\prod_{j=1}^{N-1}j!$ be the relevant Barnes product.

\begin{lemma}[HCIZ determinant inequality]\label{lem:HCIZ}
For increasingly ordered real vectors $u,v$,
\begin{equation}\label{eq:HCIZ-bound}
 0<\det(e^{u_i v_j})
 \le
 \frac{\DeltaV(u)\DeltaV(v)}{G(N+1)}
 \exp\!\left(\sum_{i=1}^N u_i v_i\right).
\end{equation}
\end{lemma}
\begin{proof}
The Harish--Chandra--Itzykson--Zuber formula with normalized Haar measure is
\[
 \int_{U(N)}e^{\operatorname{tr}(\operatorname{diag}(u)
 U\operatorname{diag}(v)U^*)}\,dU
 =G(N+1)\frac{\det(e^{u_i v_j})}{\DeltaV(u)\DeltaV(v)}.
\]
Von Neumann's trace inequality bounds the exponent in the integral above by
$\sum_i u_i v_i$.
\end{proof}

The factor $G(N+1)$ is crucial.  When both node sets come from two-dimensional
lattice regions, the leading $N^2\log N$ contributions of the two Vandermonde
products cancel the $N^2\log N$ contribution of $G(N+1)$ exactly.  The problem
therefore lives at the next, $N^2$-constant, scale.  This is the determinant
manifestation of the critical nature of four exponentials.

\subsection{Finite-prime lower bound as an assignment problem}

For a prime $p$, the valuation of an entry in~\eqref{eq:grid-matrix} is
\begin{equation}\label{eq:valuation-cost}
 w_p((a,b),(n,k))
 =\mathbf1_{p=2}an+\mathbf1_{p=3}bn+a_p ak+b_p bk.
\end{equation}

\begin{lemma}[Tropical determinant bound]\label{lem:tropical}
For every prime $p$,
\[
 v_p(D(\mathcal R,\mathcal C))
 \ge \min_{\sigma\in S_N}\sum_{i=1}^N w_p(i,\sigma(i)).
\]
\end{lemma}
\begin{proof}
Each Leibniz term has the displayed summed valuation.  The nonarchimedean
triangle inequality says that the valuation of their signed sum is at least
the minimum of their valuations.
\end{proof}

This converts divisibility into a finite optimal-transport problem.  Even a
coarse componentwise use of it already produces an $N^2$-scale divisor.

\section{Square grids: a universal divisor and the critical cancellation}

Take $\mathcal R=\mathcal C=I_m:=\{0,\ldots,m-1\}^2$ and $N=m^2$.  Denote the
corresponding determinant by $D_m$.

Let $L_m$ be the minimum of $\sum_i a_i n_{\sigma(i)}$ over bijections between
the two copies of $I_m$.  The multiset of first coordinates consists of each
integer $0,\ldots,m-1$ repeated $m$ times.  The rearrangement inequality gives
\begin{equation}\label{eq:Lm}
 L_m=m\sum_{j=0}^{m-1}j(m-1-j)
 =\frac{m^2(m-1)(m-2)}6.
\end{equation}
The same value applies to each of the four pairings $an,bn,ak,bk$.

\begin{theorem}[Universal square-grid divisor]\label{thm:square-divisor}
For every $m\ge1$,
\begin{equation}\label{eq:square-lower}
 \log D_m\ge L_m(\log2+\log3+\log A+\log B)
 =L_m(1+\beta)\log6.
\end{equation}
Equivalently, the integer $D_m$ is at least
$6^{(1+\beta)L_m}$ in logarithmic size.
\end{theorem}
\begin{proof}
For a fixed prime, the minimum of a sum of nonnegative cost components is at
least the sum of their separate minima.  Apply
Lemma~\ref{lem:tropical} to~\eqref{eq:valuation-cost}; each present component
has minimum $L_m$.  Sum the resulting lower bounds with weights $\log p$ and
use
$\sum_pa_p\log p=\log A$ and
$\sum_pb_p\log p=\log B$.
\end{proof}

The estimate is much stronger than the bare fact $D_m\ge1$.  Nevertheless, a
balanced rectangular asymptotic does not yet contradict it.  Let
$z=\sqrt{\alpha\gamma\beta}$.  After choosing rectangle aspect ratios so that
the two contributions to each projected logarithm are equal, the normalized
node distribution is that of
$Z=(X+Y)/2$ for independent $X,Y\sim\mathrm{Unif}[0,1]$.  One has
\[
 \E Z^2=\frac7{24},\qquad
 \iint\log|s-t|\,d\Law(Z)(s)d\Law(Z)(t)
 =\frac13\log2-\frac{25}{12}.
\]
Combining these with Lemma~\ref{lem:HCIZ} and
Theorem~\ref{thm:square-divisor} yields the necessary rectangular inequality
\begin{equation}\label{eq:rectangle-gap}
 h(z):=\frac12\log(4z)+\frac13\log2-\frac43+\frac z2\ge0.
\end{equation}
For the range allowed even by $A\ge3$, this inequality is already positive.
The square-grid divisor is therefore real progress in the arithmetic estimate,
but rectangular geometry spends too much archimedean size.

\section{Triangular grids and a new exclusion inequality}

The natural way to reduce the archimedean cost is to take the first $N$ nodes
rather than a rectangular block.

Let $\mathcal R_N$ consist of the $N$ pairs $(a,b)\in\Nzero^2$ with the
smallest values of $u=a\alpha+b\gamma$, and let $\mathcal C_N$ consist of the
$N$ pairs $(n,k)\in\Nzero^2$ with the smallest values of $v=n+k\beta$.
Ties do not occur.  Let $D_N$ be the positive integer determinant from these
sets.

Define
\[
 U:=\sqrt{2\alpha\gamma},\qquad
 V:=\sqrt{2\beta},\qquad
 z:=\sqrt{\alpha\gamma\beta},
\]
so $UV=2z$.

\subsection{Limiting distributions}

Elementary lattice-point counting gives
\[
 \#\{(a,b)\in\Nzero^2:a\alpha+b\gamma\le T\}
 =\frac{T^2}{2\alpha\gamma}+O(T),
\]
and similarly for $n+k\beta$.  Thus the $N$th thresholds are asymptotic to
$U\sqrt N$ and $V\sqrt N$.  After division by $\sqrt N$, the projected nodes
converge to the distributions
\[
 d\mu_U(t)=\frac{2t}{U^2}\,dt\quad(0\le t\le U),\qquad
 d\nu_V(t)=\frac{2t}{V^2}\,dt\quad(0\le t\le V).
\]
Their quantile functions are
\[
 Q_{\mu_U}(s)=U\sqrt s,\qquad Q_{\nu_V}(s)=V\sqrt s.
\]
The logarithmic energies are
\begin{equation}\label{eq:beta2-energy}
 \mathcal E(\mu_U)=\log U-\frac74,\qquad
 \mathcal E(\nu_V)=\log V-\frac74.
\end{equation}
The finite irrationality type supplied by linear forms in logarithms ensures
that exceptionally close projected lattice points do not spoil convergence of
the logarithmic energies; details are given in
Appendix~\ref{app:energy-convergence}.

\subsection{The archimedean constant}

The Barnes asymptotic is
\[
 \log G(N+1)=\frac{N^2}{2}\log N-\frac{3N^2}{4}+o(N^2).
\]
Using~\eqref{eq:beta2-energy}, the two Vandermonde terms in
Lemma~\ref{lem:HCIZ}, and comonotone coupling for the trace term, gives
\begin{equation}\label{eq:triangle-upper}
 \limsup_{N\to\infty}\frac{\log D_N}{N^2}
 \le \frac12\log(UV)-1+\frac{UV}{2}.
\end{equation}
Indeed,
\[
 \int_0^1Q_{\mu_U}(s)Q_{\nu_V}(s)\,ds
 =UV\int_0^1s\,ds=\frac{UV}{2}.
\]

\subsection{The finite-prime constant}

Uniform measure on the row triangle can be written as
\[
 a=\frac U\alpha X,\qquad b=\frac U\gamma Y,
 \qquad X,Y\ge0,\quad X+Y\le1,
\]
where $(X,Y)$ is uniform on the standard simplex.  Each coordinate has density
$2(1-x)$ and quantile
\[
 q(s)=1-\sqrt{1-s}.
\]
The column coordinates have the corresponding scales
$n=V X'$ and $k=(V/\beta)Y'$.
The countermonotone assignment constant is
\begin{equation}\label{eq:J}
 J:=\int_0^1q(s)q(1-s)\,ds
 =\frac\pi8-\frac13.
\end{equation}
Consequently each of the four logarithmically weighted components
$an\log2$, $bn\log3$, $ak\log A$, and $bk\log B$ contributes asymptotically
$UVJ N^2$ to the componentwise valuation lower bound.  Therefore
\begin{equation}\label{eq:triangle-lower}
 \liminf_{N\to\infty}\frac{\log D_N}{N^2}\ge4UVJ.
\end{equation}

\begin{theorem}[Triangular determinant inequality]\label{thm:triangle-inequality}
Every nonintegral counterexample $\beta$ satisfies
\begin{equation}\label{eq:g-ineq}
 \boxed{
 g\!\left(\sqrt{\beta\log2\log3}\right)\ge0,
 \qquad
 g(z)=\frac12\log(2z)-1+\left(\frac{11}{3}-\pi\right)z.}
\end{equation}
\end{theorem}
\begin{proof}
Subtract the lower bound~\eqref{eq:triangle-lower} from the upper
bound~\eqref{eq:triangle-upper}.  Since $UV=2z$ and
$4J=\pi/2-4/3$, the resulting necessary inequality is exactly
\eqref{eq:g-ineq}.
\end{proof}

The derivative $g'(z)=1/(2z)+11/3-\pi$ is positive for $z>0$, so $g$ has a
unique positive zero
\[
 z_0=1.12895010749713,
 \qquad
 \beta_\triangle:=\frac{z_0^2}{\log2\log3}=1.67370758736671.
\]
At the first arithmetically possible nonintegral value
$\beta=\log_2 3$, one has $z=\log3$ and
\[
 g(\log3)=-0.029549732685145<0.
\]

\begin{corollary}[A uniform small-counterexample exclusion]\label{cor:A-ge-5}
For the least nonintegral solution $\beta$,
\[
 2^\beta\ge5,
 \qquad
 \beta\ge\log_2 5=2.32192809488736,
 \qquad
 3^\beta\ge13.
\]
\end{corollary}
\begin{proof}
Theorem~\ref{thm:triangle-inequality} gives
$\beta\ge\beta_\triangle>\log_2 3$, so the integer $A=2^\beta$ is at least
$4$.  If $A=4$, then $\beta=2$ is integral.  Hence a nonintegral least
counterexample has $A\ge5$.  The final bound follows from monotonicity and
$3^{\log_2 5}=12.818619193123>12$.
\end{proof}

\begin{remark}
For a fixed small $A$, direct certified interval arithmetic can of course show
that $A^{\log_2 3}$ lies strictly between consecutive integers.  The point of
Corollary~\ref{cor:A-ge-5} is different: it is the first output of a uniform
asymptotic determinant inequality, and its derivation exposes the arithmetic
constant that must be improved to attack all $A$ simultaneously.
\end{remark}

\subsection{Numerical anatomy of the two grid geometries}

The table compares the triangular gap $g$ with the balanced-rectangle gap $h$
from~\eqref{eq:rectangle-gap}.  Negative values would contradict a
counterexample at that $\beta$.

\begin{center}
\begin{tabular}{lrrr}
\toprule
$\beta$ & $z=\sqrt{\alpha\gamma\beta}$ & $g(z)$ & $h(z)$ \\
\midrule
$1$ & 0.8726396796 & -0.2633422644 & -0.04093352603 \\
$\log_2 3$ & 1.098612289 & -0.02954973269 & 0.1871929656 \\
$\beta_{\triangle}$ & 1.128950107 & -1.054219794e-81 & 0.2159820074 \\
$2$ & 1.23409887 & 0.09973735836 & 0.3130828643 \\
$\log_2 5$ & 1.329717364 & 0.1872568287 & 0.3982047952 \\
$3$ & 1.511456262 & 0.3467367942 & 0.5531278372 \\
$5$ & 1.951281643 & 0.7053840783 & 0.900746934 \\
$10$ & 2.759528964 & 1.303060538 & 1.47815739 \\
\bottomrule
\end{tabular}
\end{center}

The triangular region improves the constant substantially, but $g$ is
increasing and becomes comfortably positive in the range $\beta\ge\log_2 5$.
To continue, one must either improve the finite-prime lower bound, find a
better variational region, or enrich the determinant.

\section{Rational changes of basis and the product formula}

A natural reaction to the positive gap is to replace $3$ by a ratio such as
$3/2$, whose logarithm is smaller.  Ignoring denominators makes this look like
an immediate proof.  The product formula shows exactly why that argument is
false and, more usefully, turns it into quantitative denominator constraints.

\subsection{A signed arithmetic functional}

For a rational $q>1$, write it in lowest terms as
$q=\num(q)/\den(q)$ and define
\begin{equation}\label{eq:rho}
 \rho(q):=
 \frac{J\log\num(q)-K\log\den(q)}{\log q},
 \qquad
 J=\frac\pi8-\frac13,\quad K=\frac16.
\end{equation}
The constant $J$ is the countermonotone coordinate cost in a simplex, while
$K$ is the comonotone cost, since a simplex coordinate has second moment
$1/6$.  If $q$ is an integer, then $\rho(q)=J$.  A denominator contributes with
$K$, not $J$, because a negative valuation reverses the relevant assignment
bound.

\begin{theorem}[Adelic triangular inequality]\label{thm:adelic-triangle}
Let $r_1,r_2>1$ be multiplicatively independent positive rationals and suppose
$R_i=r_i^\beta\in\Q_{>1}$ for $i=1,2$.  Put
\[
 z_r:=\sqrt{\beta\log r_1\log r_2}.
\]
Then the existence of these four rational exponentials implies
\begin{equation}\label{eq:adelic-rho}
 \frac12\log(2z_r)-1+z_r
 -2z_r\bigl(\rho(r_1)+\rho(r_2)+\rho(R_1)+\rho(R_2)\bigr)\ge0.
\end{equation}
Equivalently, if
\[
 \Delta_{\rm den}:=
 \sum_{q\in\{r_1,r_2,R_1,R_2\}}
 \frac{\log\den(q)}{\log q},
\]
then
\begin{equation}\label{eq:adelic-den}
 g(z_r)+2z_r\left(\frac12-\frac\pi8\right)\Delta_{\rm den}\ge0.
\end{equation}
\end{theorem}
\begin{proof}
Use the same first-$N$ triangular sets for the positive logarithms
$\log r_1,\log r_2$ and for $1,\beta$.  The determinant entries are now
rational.  For a fixed prime, split each signed valuation coefficient into its
positive and negative parts.  In the tropical determinant bound, a positive
coefficient is multiplied by the countermonotone minimum $J$, while a negative
coefficient is multiplied by the comonotone maximum $K$.  Summing over primes
and using the product formula for the nonzero rational determinant gives the
lower constant
\[
 2z_r\bigl(\rho(r_1)+\rho(r_2)+\rho(R_1)+\rho(R_2)\bigr).
\]
The HCIZ upper constant remains $\tfrac12\log(2z_r)-1+z_r$.  This proves
\eqref{eq:adelic-rho}.  Finally,
\[
 \rho(q)=J-(K-J)\frac{\log\den(q)}{\log q},
\]
and substitution gives~\eqref{eq:adelic-den}.
\end{proof}

\subsection{All unimodular bases of $\langle2,3\rangle$}

Let
\[
 M=\begin{pmatrix}a&b\\c&d\end{pmatrix}\in\operatorname{GL}_2(\Z)
\]
and suppose
\[
 \lambda_1=a\alpha+b\gamma>0,\qquad
 \lambda_2=c\alpha+d\gamma>0.
\]
Set
\[
 r_1=2^a3^b,\quad r_2=2^c3^d,\qquad
 R_1=A^aB^b,\quad R_2=A^cB^d.
\]
Then Theorem~\ref{thm:adelic-triangle} applies.  This supplies an infinite
family of necessary inequalities indexed by unimodular changes of basis.
Continued-fraction vectors make one $\lambda_i$ small, but the corresponding
rational base has a large denominator; Equation~\eqref{eq:adelic-den} measures
the precise price.

\begin{proposition}[Optimality of the integral basis]\label{prop:integral-basis-opt}
Among unimodular bases for which both transformed bases are integers, the
product $\lambda_1\lambda_2$ is at least $\alpha\gamma$.  Equality occurs only
for the original basis $(2,3)$ up to permutation.  Likewise, among unimodular
bases of the solution lattice whose two transformed exponents remain in
$\mathcal S$, the product of the positive exponent generators is at least
$\beta$.
\end{proposition}
\begin{proof}
If $2^a3^b$ is an integer, then $a,b\ge0$.  Thus an integral unimodular basis
has nonnegative entries.  Expanding,
\[
 (a\alpha+b\gamma)(c\alpha+d\gamma)
 \ge(ad+bc)\alpha\gamma\ge\alpha\gamma,
\]
because $|ad-bc|=1$.  Equality forces the standard basis up to permutation.
For the exponent lattice, Theorem~\ref{thm:exact-monoid} says that a transformed
exponent is simultaneously integral at bases $2$ and $3$ exactly when its
coefficients in $(1,\beta)$ are nonnegative.  The same argument applies.
\end{proof}

This proposition explains why a denominator-free lattice reduction cannot
improve the triangular inequality: the exact monoid has already selected the
smallest integral basis.

\subsection{The basis $(2,3/2)$ and a gcd constraint}

Take
\[
 r_1=2,\qquad r_2=\frac32,\qquad
 R_1=A,\qquad R_2=\frac BA.
\]
Let $d=A/\gcd(A,B)$ be the denominator of $B/A$ in lowest terms and put
$\delta=\log(3/2)$.  Then
\begin{equation}\label{eq:delta-32}
 \Delta_{\rm den}=rac{\alpha}{\delta}
 +\frac{\log d}{\beta\delta}.
\end{equation}
Theorem~\ref{thm:adelic-triangle} therefore gives a rigorous lower bound on
$d$, whenever the right side obtained by solving~\eqref{eq:adelic-den} is
positive.  In words: if $3/2$ substantially improves the archimedean
configuration, then $B/A$ must carry enough denominator to pay for that
improvement.  A counterexample cannot have excessive divisibility
$A\mid B$ in the range where the reduced-basis gap is negative.

For orientation, the next table substitutes the hypothetical value
$\beta=\log_2 A$.  The column $d_{\min}$ is the least integer denominator not
excluded by the inequality alone; it is not a claim that the corresponding
$B$ exists.

\begin{center}
\begin{longtable}{rrrrr}
\toprule
$A$ & $\log_2A$ & $z_r$ & gap before $\log d$ & $d_{\min}$ \\
\midrule
\endfirsthead
\toprule
$A$ & $\log_2A$ & $z_r$ & gap before $\log d$ & $d_{\min}$ \\
\midrule
\endhead
3 & 1.5849625 & 0.667419621 & -0.260297461 & 4 \\
5 & 2.32192809 & 0.807818616 & -0.0396109308 & 2 \\
6 & 2.5849625 & 0.852347316 & 0.0269341919 & 1 \\
7 & 2.80735492 & 0.88825597 & 0.0795954278 & 1 \\
9 & 3.169925 & 0.94387388 & 0.159569469 & 1 \\
10 & 3.32192809 & 0.966239056 & 0.191227186 & 1 \\
11 & 3.45943162 & 0.986033907 & 0.219022705 & 1 \\
12 & 3.5849625 & 1.00376439 & 0.243748126 & 1 \\
13 & 3.70043972 & 1.01980266 & 0.265979188 & 1 \\
14 & 3.80735492 & 1.03443012 & 0.286146736 & 1 \\
15 & 3.9068906 & 1.04786443 & 0.304581094 & 1 \\
17 & 4.08746284 & 1.07180649 & 0.337231592 & 1 \\
18 & 4.169925 & 1.08256404 & 0.351820061 & 1 \\
19 & 4.24792751 & 1.09264233 & 0.36544257 & 1 \\
20 & 4.32192809 & 1.10211837 & 0.378212214 & 1 \\
21 & 4.39231742 & 1.11105698 & 0.39022374 & 1 \\
22 & 4.45943162 & 1.11951323 & 0.401557272 & 1 \\
23 & 4.52356196 & 1.12753426 & 0.412281139 & 1 \\
\bottomrule
\end{longtable}
\end{center}

A second useful basis is
\[
 r_1=\frac32,\qquad r_2=\frac43,
 \qquad R_1=\frac BA,\qquad R_2=\frac{A^2}{B}.
\]
Its two logarithms are smaller, but Equation~\eqref{eq:adelic-den} now charges
the denominators of all four displayed ratios.  It therefore constrains the
possibility that $B/A$ and $A^2/B$ are simultaneously close to integers.  More
generally, continued-fraction bases produce a network of constraints on
\[
 \frac{B^q}{A^p}
 \quad\text{when}\quad q\log3-p\log2\text{ is small}.
\]
This is a concrete future route: prove that the denominator requirements for a
carefully chosen finite family of bases are incompatible with the prime
valuation vectors of $A$ and $B$.

\section{Joint prime transport and shared-factor gains}

The componentwise lower bounds used above are intentionally conservative.  At
a fixed prime, a single permutation must handle all valuation components at
once.  If a prime occurs in more than one generator, the separate
countermonotone assignments can conflict, making the true minimum strictly
larger.

For the original triangular determinant, define
\[
 \mathcal T_p(N):=\frac1{N^2}
 \min_{\sigma\in S_N}\sum_i w_p(i,\sigma(i)).
\]
Then the strongest direct product-formula lower bound is
\[
 \liminf_{N\to\infty}\frac{\log D_N}{N^2}
 \ge\sum_p(\log p)\liminf_{N\to\infty}\mathcal T_p(N).
\]
The baseline $4UVJ$ replaces each $\mathcal T_p$ by the sum of four independent
one-coordinate minima.

A simple rigorous gain estimate is available for a prime $p\notin\{2,3\}$ that
divides both $A$ and $B$.  Put $r=v_p(A)$, $s=v_p(B)$ and
\[
 R=\frac r\alpha,\qquad S=\frac s\gamma.
\]
On the normalized row simplex, the relevant scalar is $RX+SY$, while the
column coordinate $k$ has the marginal law of a simplex coordinate.  The
countermonotone cost is at least the product of the means minus the product of
the standard deviations.  Since
\[
 \E X=\E Y=\frac13,\quad
 \Var X=\Var Y=\frac1{18},\quad
 \Cov(X,Y)=-\frac1{36},
\]
one obtains
\begin{equation}\label{eq:overlap-bound}
 \liminf\mathcal T_p(N)
 \ge\frac{UV}{\beta}
 \left(
  \frac{R+S}{9}
  -\frac{\sqrt{R^2+S^2-RS}}{18}
 \right).
\end{equation}
The componentwise baseline for the same two terms is
\[
 \frac{UV}{\beta}J(R+S).
\]
Taking the larger of these two lower bounds gives a fully explicit shared-prime
improvement.  Summed over common prime divisors of $A$ and $B$, it yields a
sufficient contradiction criterion whenever the gain exceeds the positive
gap $g(z)$.

The estimate is strongest when $r/\alpha$ and $s/\gamma$ are comparable and
vanishes continuously as one valuation goes to zero.  It therefore does not
supply a universal proof: the worst arithmetic configuration may have nearly
disjoint prime support.  But it turns ``a large gcd should help'' into a
quantitative theorem and identifies exact valuation patterns worth targeting.

\section{Universal congruence divisibility is lower order}

There is also divisibility at primes not dividing $6AB$.  Modulo such a prime
$p$, a row of the matrix depends on $(a,b)$ only modulo $p-1$, by Fermat's
little theorem.  Hence there are at most $(p-1)^2$ distinct rows modulo $p$.
For an $N\times N$ determinant,
\[
 v_p(D_N)\ge N-\operatorname{rank}_{\mathbb F_p}(M)
 \ge N-(p-1)^2
\]
whenever the right side is positive.  Summing over primes $p\le\sqrt N$ and
using the prime number theorem gives the universal extra estimate
\begin{equation}\label{eq:congruence-gain}
 \sum_{\substack{p\le\sqrt N\\p\nmid6AB}}
  \bigl(N-(p-1)^2\bigr)\log p
 =\left(\frac23+o(1)\right)N^{3/2}.
\end{equation}
Higher prime powers and column congruences can improve the lower-order term.
They do not alter the $N^2$ constant in
Theorem~\ref{thm:triangle-inequality}.  This provides a useful barrier result:
congruence collisions that depend only on finite-group periodicity are too
small by a factor of $\sqrt N$.  A complete determinant proof needs
counterexample-specific valuation structure or a different determinant
geometry.

\section{Why several tempting routes still stop short}

\subsection{Continued fractions and integer gaps}

For integers $p,q>0$,
\[
 \frac{A^q}{2^p}=2^{q\beta-p},\qquad
 \frac{B^q}{3^p}=3^{q\beta-p}.
\]
When $p/q$ is a convergent to $\beta$, both ratios are close to $1$.  Since the
cross-multiplied differences are nonzero integers, one obtains exponential
lower bounds such as
\[
 |q\beta-p|\gg 2^{-p}.
\]
These are far weaker than the polynomial lower bounds already supplied by
linear forms in logarithms.  Conversely, Baker's polynomial lower bounds do
not contradict the $q^{-2}$ approximations supplied by continued fractions.
The two representations of $\beta$ give proportional linear forms, not two
independent small quantities.

\subsection{Catalan, $abc$, and small differences}

The equations $A^q-2^p=\pm1$ and $B^q-3^p=\pm1$ would be highly constrained,
but continued fractions only make the relative difference small, not the
absolute difference.  The absolute scale is exponential in $q$.  Standard
$abc$ heuristics do not bridge that exponential gap without an additional
source of unusually good approximation.

\subsection{Counting rational points}

The orbit of fractional parts gives $O(\log H)$ rational points of height at
most $H$ on a fixed transcendental arc.  This is compatible with every known
subpolynomial upper bound for rational points on definable transcendental
curves.  A counting theorem would need to exploit the pure $2$-power and
$3$-power denominator restrictions, not merely rationality.

\subsection{The four-exponentials boundary}

If $x$ is irrational, the pairs $(\log2,\log3)$ and $(1,x)$ are both
$\Q$-linearly independent, while all four exponentials
\[
 2,\quad3,\quad2^x,\quad3^x
\]
are algebraic.  Thus the Four Exponentials Conjecture would prove
Conjecture~\ref{conj:AE} immediately.  The exact cancellation of the
$N^2\log N$ terms in our determinant estimates is the quantitative shadow of
being exactly at this boundary.  Six Exponentials is sufficient to collapse
all possible exceptions to one affine direction, but not to remove that last
direction.

\section{A variational reformulation of the remaining determinant problem}

The preceding calculations extend to scaled lattice regions
$R,C\subset\R_{\ge0}^2$ of equal area.  Push forward uniform measure on $R$ by
$(a,b)\mapsto a\alpha+b\gamma$, and uniform measure on $C$ by
$(n,k)\mapsto n+k\beta$.  The HCIZ side is determined by their logarithmic
energies and comonotone transport.  The finite-prime side is a sum of
countermonotone or, more accurately, joint optimal-transport costs for the
valuation bilinear forms~\eqref{eq:valuation-cost}.

This suggests the following concrete variational program.
\begin{enumerate}
\item Choose regions $R,C$ and compute the archimedean functional
\[
 \mathcal A(R,C)=\frac12\mathcal E(\mu_R)
 +\frac12\mathcal E(\nu_C)+\frac34
 +\int_0^1Q_{\mu_R}(t)Q_{\nu_C}(t)\,dt,
\]
with the scale normalization determined by $\operatorname{area}(R)
=\operatorname{area}(C)$.
\item Compute or bound the exact prime transport functional
\[
 \mathcal P(R,C;A,B)=
 \sum_p\log p\inf_{\pi}
 \int w_p(r,\pi(r))\,dr.
\]
\item Any counterexample must satisfy
$\mathcal A(R,C)\ge\mathcal P(R,C;A,B)$ for every admissible pair of regions.
A single strict reverse inequality proves a contradiction.
\end{enumerate}

Rectangles and origin triangles are analytically soluble test cases.  The next
step is not merely to scan aspect ratios: it is to optimize jointly over region
shape and valuation configurations allowed by
Corollary~\ref{cor:unit} and Lemma~\ref{lem:primitive-gcd}.  The latter
constraints are discrete and may prevent the valuation directions that make
the transport cost artificially small.

\section{Confluent and integer-valued determinant extensions}

Ordinary evaluations use one row per lattice point.  A potentially stronger
construction uses integer jets.  View~\eqref{eq:grid-matrix} as evaluating the
monomials $X^nY^k$ at the integer points
\[
 (X,Y)=(2^a3^b,A^aB^b).
\]
The Euler operators
\[
 D_X=X\frac{\partial}{\partial X},\qquad
 D_Y=Y\frac{\partial}{\partial Y}
\]
act by
\[
 D_X^rD_Y^s(X^nY^k)=n^rk^sX^nY^k,
\]
which remains integral at the lattice points.  Hasse derivatives give the
alternative integer values
\[
 \binom nr\binom ksX^{n-r}Y^{k-s}.
\]
Thus one can build confluent evaluation determinants without introducing the
irrational factors $n+k\beta$ that ordinary derivatives along the real curve
would produce.

The unresolved issues are nonvanishing and optimal scaling.  A useful target
is an integer-jet determinant whose archimedean limit behaves like a confluent
HCIZ determinant while its $p$-adic assignment problem acquires an additional
$N^2$ contribution.  This is more promising than generic factorial
normalization: the elementary divisibility obtained from binomial bases alone
is only of order $N^{3/2}\log N$ for balanced two-dimensional grids.

\section{Conditional resolutions and exact reformulations}

\subsection{Four Exponentials and Schanuel}

The Four Exponentials Conjecture immediately implies
Conjecture~\ref{conj:AE}.  Schanuel's conjecture implies Four Exponentials and
therefore also suffices.  The present problem is thus a particularly concrete
integral special case of a central algebraic-independence conjecture.

\subsection{Logarithm determinant formulation}

With $A=2^x$ and $B=3^x$, the condition is
\[
 \det\begin{pmatrix}
 \log2&\log A\\
 \log3&\log B
 \end{pmatrix}=0.
\]
The conjecture says that if $A,B$ are positive integers and this determinant
vanishes, then the two columns are related by a nonnegative integer factor.
Baker's theorem controls nonzero \emph{linear} forms in logarithms, whereas
this is a bilinear relation among four logarithms; that distinction is exactly
where standard transcendence machinery stops.

\subsection{Order-preserving monoid formulation}

A counterexample would give an order-preserving monoid isomorphism
\[
 \langle2,A\rangle_{\Nzero^2}\longrightarrow
 \langle3,B\rangle_{\Nzero^2},
 \qquad2^nA^k\mapsto3^nB^k,
\]
because both sides have logarithm $n+k\beta$ up to a fixed positive scale.
Theorem~\ref{thm:exact-monoid} says this monoid is the entire domain on which
$m\mapsto m^{\log_2 3}$ is integer-valued.  A purely arithmetic classification
of such order-preserving power isomorphisms would prove the conjecture without
explicit transcendence estimates.

\section{Formalization roadmap}

A substantial portion of the new work is suitable for formalization in Lean,
with the deep analytic inputs isolated behind explicit interfaces.

\begin{enumerate}
\item Define $\mathcal S$ and prove the range, rationality, additive-monoid, and
      local-finiteness lemmas using elementary real analysis and unique
      factorization.
\item Package the Six Exponentials Theorem as an external theorem interface.
      From it, formalize Proposition~\ref{prop:affine-rank}.
\item Use finitely supported prime-valuation maps for $A$ and $B$ to formalize
      Lemmas~\ref{lem:primitive-gcd} and
      \ref{lem:denominator-elimination}.  This is the most valuable near-term
      target: it converts an external transcendence theorem into the exact
      monoid identity using only discrete arithmetic.
\item Define finite matrices~\eqref{eq:grid-matrix}.  The tropical valuation
      lemma follows directly from the Leibniz formula and valuation of a sum.
      The square-grid assignment value~\eqref{eq:Lm} is a finite rearrangement
      theorem.
\item Treat strict total positivity, HCIZ, logarithmic-energy convergence, and
      asymptotic transport as separate analytic interfaces.  Finite versions
      of all determinant inequalities can still be machine checked once
      supplied with rational interval certificates.
\item For computational searches over region shapes or rational bases, emit
      proof certificates consisting of finite node lists, exact valuation
      assignment dual variables, and rational upper bounds for the real
      determinant.
\end{enumerate}

This separation would let the ProveIt development verify every new arithmetic
step even before the transcendence and asymptotic libraries are available.

\section{Prioritized research directions}

\begin{enumerate}
\item \textbf{Joint valuation optimization.}  Replace the componentwise
assignment lower bound by the exact transport cost for each prime, and optimize
over valuation vectors satisfying the primitive-gcd and base-prime-unit
conditions.  This is the most direct way to seek the missing $N^2$ constant.

\item \textbf{Rational-basis denominator network.}  Apply
Theorem~\ref{thm:adelic-triangle} to a finite set of continued-fraction bases.
Translate the resulting inequalities into lower bounds on denominators of
$B^q/A^p$.  Try to show that unique factorization makes the bounds mutually
incompatible.

\item \textbf{Variational region search.}  Numerically optimize convex and
nonconvex lattice regions using the exact archimedean energy and assignment
functionals.  Any promising negative gap should then be converted into a
rigorous interval proof.

\item \textbf{Integer jets and confluent total positivity.}  Establish a
nonvanishing theorem for Euler/Hasse-jet determinants and determine whether
their prime transport grows faster than their archimedean capacity.

\item \textbf{Exploit minimal-generator valuations.}  The conditions
$\min(v_2(A),v_3(B))=0$ and $g_0=1$ have not yet been fully inserted into the
variational optimization.  They should be imposed from the beginning rather
than used only as post-processing constraints.

\item \textbf{Effective finite certificates.}  For candidate ranges of $A$,
combine rigorous interval arithmetic with determinant and gcd constraints.
Although finite verification cannot settle the conjecture alone, it can reveal
which valuation patterns come closest to saturating the theoretical bounds.
\end{enumerate}


\section{A proof-producing finite verification}
\label{sec:finite-verification}

The structural and determinant arguments above are uniform, but they leave an
infinite range of possible values of $A=2^\beta$.  A complementary finite
calculation can remove an initial segment without assuming any unproved
transcendence statement.  The following computation is deliberately based on
outward-rounded intervals and an explicit positive-series remainder, rather
than on a test such as ``the floating-point value is not visually close to an
integer.''

\begin{theorem}[Certified finite range]
\label{thm:finite-range}
For every integer $A$ with
\[
  2\le A\le 100,000,
\]
the number $A^{\log 3/\log 2}$ is an integer if and only if $A$ is a power
of $2$.  Consequently, a nonintegral Alaoglu--Erd\H{o}s counterexample must
satisfy
\[
  2^\beta>100,000,
  \qquad
  \beta>\log_2(100,000)\approx 16.609640474437.
\]
\end{theorem}

\begin{proof}[Certificate method]
For fixed $A$, put
\[
  F_A(b)=(\log b)(\log 2)-(\log A)(\log 3),\qquad b>0.
\]
The function $F_A$ is strictly increasing, and
$A^{\log 3/\log 2}=b$ is equivalent to $F_A(b)=0$.
Every logarithm is enclosed without calling a transcendental floating-point
routine in the decisive comparison.  For $0\le z\le 1/3$ we use
\[
  \log\frac{1+z}{1-z}
  =2\sum_{j=0}^{M-1}\frac{z^{2j+1}}{2j+1}+R_M(z),
  \qquad
  0<R_M(z)<
  \frac{2z^{2M+1}}{(2M+1)(1-z^2)}.
\]
For an integer $n\ge1$, choose $k=\lfloor\log_2 n\rfloor$ and write
\[
  \log n=k\log2+
  \log\frac{1+z_n}{1-z_n},
  \qquad
  z_n=\frac{n-2^k}{n+2^k}\in[0,1/3).
\]
Thus all terms are positive rational operations, and directed rounding gives
certified lower and upper endpoints.  The run used 70 decimal digits and
$M=28$ terms.  A binary64 value was used only as an initial guess for the
integer immediately below the root; the program then moved monotonically
until it had certified, for consecutive integers $b,b+1$,
\[
  F_A(b)<0<F_A(b+1).
\]
Hence no integer can be the root.  Powers of $2$ are separated at the start,
and for $A=2^m$ the corresponding value is exactly $3^m$.

The completed run checked 99,983 nonpowers of $2$.  The largest number of
integer steps needed to repair an initial binary64 guess was 0.  The
smallest absolute interval endpoint encountered was
$\texttt{6.263278510736485942331545732E-13}$, far outside the interval uncertainty.  Runtime on the
working container was approximately 15.238 seconds.  The exact verifier
used for the run is reproduced in Appendix~\ref{app:finite-verifier}; its
SHA--256 digest is
\begin{center}
\texttt{acc099fb776f830b243a2421bc7ce27d5bedc63a38c3a7b190accb3c621cbadd}.
\end{center}
\end{proof}

\begin{remark}
The numerical bound is not presented as evidence that near misses cannot occur
later.  Its role is narrower and rigorous: it turns all $A$ in the displayed
finite interval into closed cases, while the determinant inequalities address
whole infinite parameter ranges and arithmetic configurations at once.
\end{remark}

\section{Conclusion}

The continuation does not close the final four-exponentials gap, but it
substantially narrows and clarifies the problem.

First, the hypothetical exceptional set is completely rigid: one least
counterexample $\beta$ would generate every solution, and only solutions, as
$n+k\beta$ with $n,k\ge0$.  Second, that exact grid supports determinants that
are simultaneously strictly totally positive over the reals and highly
divisible at finite primes.  The square-grid divisor and the triangular
inequality quantify both sides at the correct $N^2$ scale.  Third, rational
basis reduction is no longer a vague temptation: the adelic inequality shows
exactly how denominators pay for every smaller logarithm.  Finally, the
remaining obstacle has a precise form---one needs an additional $N^2$-scale
arithmetic gain, not merely more lattice points or a sharper polynomial error
term.

The most promising immediate attack is therefore a joint one: optimize the
triangular or variational determinant using the full prime valuation vectors,
while applying several rational bases to force incompatible denominator and
gcd requirements.  That program is concrete, computationally testable, and
amenable to partial formal verification.

\appendix

\section{Logarithmic-energy convergence for projected lattice triangles}
\label{app:energy-convergence}

We record the technical point needed in the triangular determinant limit.  Let
$\lambda_1,\lambda_2>0$ and assume their ratio has finite irrationality type:
there are $c,\tau>0$ such that
\[
 |r\lambda_1+s\lambda_2|\ge c(1+|r|+|s|)^{-\tau}
\]
for all nonzero $(r,s)\in\Z^2$.  Let $t_{1,N}<\cdots<t_{N,N}$ be the first $N$
values of $a\lambda_1+b\lambda_2$ with $a,b\in\Nzero$, and set
$\widetilde t_{i,N}=t_{i,N}/\sqrt N$.  Then the empirical measures converge to
\[
 d\mu(t)=\frac{2t}{U^2}\,dt,\qquad U=\sqrt{2\lambda_1\lambda_2},
\]
and
\begin{equation}\label{eq:energy-limit-app}
 \frac{2}{N^2}\sum_{i<j}\log|\widetilde t_{j,N}-\widetilde t_{i,N}|
 \longrightarrow\iint\log|x-y|\,d\mu(x)d\mu(y).
\end{equation}

The weak convergence follows from a Riemann-sum lattice count in the triangle
$\lambda_1a+\lambda_2b\le U\sqrt N$.  For the singular kernel, truncate
$\log|x-y|$ below at $-M$.  Weak convergence handles the truncated kernel.
Uniform integrability of the negative tail follows by counting difference
vectors $(r,s)$ in thin strips around
$r\lambda_1+s\lambda_2=0$; the finite-type lower bound prevents an individual
difference from being smaller than a fixed power of $N$, while the strip count
shows that the proportion below $\varepsilon$ is $O(\varepsilon)$ up to a
power-saving discrepancy term.  Integrating the distribution function of
$-\log|x-y|$ and then sending $M\to\infty$ proves
\eqref{eq:energy-limit-app}.

For $(\lambda_1,\lambda_2)=(\log2,\log3)$, the required finite type follows from
linear forms in logarithms.  For $(1,\beta)$, it follows from
$\beta=\log A/\log2$ and the multiplicative independence of $A$ and $2$.

\section{Integral evaluations used in the determinant constants}

\subsection{The beta-two logarithmic energy}

For the probability density $f(x)=2x$ on $[0,1]$,
\begin{align*}
 \iint_0^1 4xy\log|x-y|\,dx\,dy
 &=8\int_0^1\int_0^xxy\log(x-y)\,dy\,dx\\
 &=8\int_0^1x^3\left(\frac12\log x-\frac34\right)dx\\
 &=-\frac74.
\end{align*}
Scaling by $U$ adds $\log U$, proving
Equation~\eqref{eq:beta2-energy}.

\subsection{The simplex assignment constant}

A coordinate of a point uniformly distributed on the standard two-simplex has
quantile $q(s)=1-\sqrt{1-s}$.  Therefore
\begin{align*}
 J&=\int_0^1q(s)q(1-s)\,ds\\
 &=\int_0^1\left(1-\sqrt s-\sqrt{1-s}+\sqrt{s(1-s)}\right)ds\\
 &=1-\frac43+\frac\pi8
 =\frac\pi8-\frac13.
\end{align*}
Its comonotone second-moment constant is
$K=\int_0^1q(s)^2ds=1/6$.

\subsection{The balanced-rectangle energy}

Let $Z=(X+Y)/2$ with independent uniform $X,Y\in[0,1]$.  The density of
$S=X+Y$ is triangular, and the density of the difference of two independent
copies of $S$ is
\[
 h(t)=
 \begin{cases}
  \frac23-t^2+\frac12t^3,&0\le t\le1,\\
  \frac16(2-t)^3,&1\le t\le2.
 \end{cases}
\]
A direct integration gives
\[
 \E\log|S-S'|=\frac43\log2-\frac{25}{12}.
\]
Since $Z-Z'=(S-S')/2$,
\[
 \E\log|Z-Z'|=\frac13\log2-\frac{25}{12}.
\]
Also $\E Z^2=7/24$.

\section{General two-prime version}

The structural and determinant arguments extend verbatim to distinct primes
$p,q$.  If a nonintegral $\beta$ satisfies $p^\beta,q^\beta\in\Z$, choose the
least such $\beta$.  Six Exponentials gives the same affine-rank collapse, and
valuation minimality yields the exact monoid $\Nzero+\Nzero\beta$.  The
triangular determinant then gives
\[
 g\!\left(\sqrt{\beta\log p\log q}\right)\ge0.
\]
This form may be useful for comparing the strength of the method across base
pairs and for testing variational improvements.

\section{Finite determinant certificates}

The asymptotic argument can be converted into finite certificates.  For chosen
finite node sets, a certificate consists of:
\begin{enumerate}
\item the exact integer or rational matrix entries;
\item sorted rational intervals for $u_i,v_i$ and their differences;
\item an upper bound from Equation~\eqref{eq:HCIZ-bound};
\item for each relevant prime, dual potentials for the assignment problem in
      Lemma~\ref{lem:tropical}, proving the claimed minimum;
\item a product-formula comparison showing the upper bound is smaller than the
      certified divisor.
\end{enumerate}
Such certificates are compact enough for independent checking and eventual
Lean verification.  They also avoid trusting floating-point determinant
computations, which are extremely ill-conditioned in the near-singular regime.


\section{Source of the finite verifier}
\label{app:finite-verifier}

The following is the complete source used to establish
Theorem~\ref{thm:finite-range}.  Its only non-rigorous arithmetic---the
binary64 call used to suggest a starting integer---cannot affect the result:
the subsequent directed-interval sign search is monotone and independently
certifies the final pair of inequalities.

\begin{lstlisting}[style=proofpython]
#!/usr/bin/env python3
"""Rigorous finite verification for the Alaoglu--Erdos problem.

For every integer A in [2, LIMIT], this proves that A^(log(3)/log(2))
is an integer only when A is a power of two.  All decisive comparisons use
Decimal directed rounding plus the positive atanh series for logarithms.
The ordinary binary64 approximation is used only to choose an initial nearby
integer; a monotone certified search repairs any inaccurate guess.
"""
from __future__ import annotations

from dataclasses import dataclass
from decimal import Decimal, Context, ROUND_FLOOR, ROUND_CEILING, localcontext
import json
import math
import sys
import time

LIMIT = int(sys.argv[1]) if len(sys.argv) > 1 else 100_000
PREC = 70
TERMS = 28

DOWN = Context(prec=PREC, rounding=ROUND_FLOOR)
UP = Context(prec=PREC, rounding=ROUND_CEILING)

@dataclass(frozen=True)
class Interval:
    lo: Decimal
    hi: Decimal


def _atanh_log_ratio(p: int, q: int, terms: int = TERMS) -> Interval:
    """Enclose 2*atanh(p/q), for 0 <= p/q <= 1/3."""
    assert 0 <= p <= q // 3 + (1 if 3 * p == q else 0)
    if p == 0:
        z = Decimal(0)
        return Interval(z, z)

    with localcontext(DOWN):
        z_lo = Decimal(p) / Decimal(q)
        z2_lo = z_lo * z_lo
        power = z_lo
        total = Decimal(0)
        for j in range(terms):
            total = total + power / Decimal(2 * j + 1)
            power = power * z2_lo
        lo = Decimal(2) * total

    with localcontext(UP):
        z_hi = Decimal(p) / Decimal(q)
        z2_hi = z_hi * z_hi
        power = z_hi
        total = Decimal(0)
        for j in range(terms):
            total = total + power / Decimal(2 * j + 1)
            power = power * z2_hi
        partial = Decimal(2) * total
        # power is now z^(2*terms+1).  Since all omitted terms are positive,
        # replace every odd denominator by the first one and sum geometrically.
        rem = Decimal(2) * power / (
            Decimal(2 * terms + 1) * (Decimal(1) - z2_hi)
        )
        hi = partial + rem
    return Interval(lo, hi)


LOG2 = _atanh_log_ratio(1, 3)


def log_integer(n: int, terms: int = TERMS) -> Interval:
    """Rigorous interval for log(n), n >= 1, by binary range reduction."""
    assert n >= 1
    if n == 1:
        z = Decimal(0)
        return Interval(z, z)
    k = n.bit_length() - 1
    two_k = 1 << k
    # n/2^k in [1,2), and z=(y-1)/(y+1) lies in [0,1/3).
    local = _atanh_log_ratio(n - two_k, n + two_k, terms)
    with localcontext(DOWN):
        lo = Decimal(k) * LOG2.lo + local.lo
    with localcontext(UP):
        hi = Decimal(k) * LOG2.hi + local.hi
    return Interval(lo, hi)


LOG3 = log_integer(3)


def determinant_log_interval(log_a: Interval, b: int) -> Interval:
    """Enclose F(A,b)=log(b)log(2)-log(A)log(3)."""
    log_b = log_integer(b)
    with localcontext(DOWN):
        lo = log_b.lo * LOG2.lo - log_a.hi * LOG3.hi
    with localcontext(UP):
        hi = log_b.hi * LOG2.hi - log_a.lo * LOG3.lo
    return Interval(lo, hi)


def certified_sign(log_a: Interval, b: int) -> tuple[int, Interval]:
    f = determinant_log_interval(log_a, b)
    if f.hi < 0:
        return -1, f
    if f.lo > 0:
        return +1, f
    # At this finite precision an interval meeting zero is not silently accepted.
    raise ArithmeticError(f"inconclusive interval at A-log={log_a}, b={b}: {f}")


def is_power_of_two(n: int) -> bool:
    return n > 0 and (n & (n - 1)) == 0


def main() -> None:
    start = time.monotonic()
    rho_approx = math.log(3.0) / math.log(2.0)
    min_abs_endpoint = Decimal('Infinity')
    max_repairs = 0
    checked = 0

    for a in range(2, LIMIT + 1):
        if is_power_of_two(a):
            continue
        checked += 1
        log_a = log_integer(a)
        guess = max(1, int(math.floor(math.exp(rho_approx * math.log(float(a))))))
        repairs = 0

        # Find consecutive integers b,b+1 on opposite sides of the unique real root.
        while True:
            s0, f0 = certified_sign(log_a, guess)
            s1, f1 = certified_sign(log_a, guess + 1)
            min_abs_endpoint = min(
                min_abs_endpoint,
                abs(f0.lo), abs(f0.hi), abs(f1.lo), abs(f1.hi)
            )
            if s0 < 0 < s1:
                break
            if s1 < 0:
                guess += 1
            elif s0 > 0:
                guess -= 1
                assert guess >= 1
            else:
                raise AssertionError((a, guess, s0, s1))
            repairs += 1
            if repairs > 100:
                raise AssertionError(f"unexpectedly bad initial approximation at A={a}")
        max_repairs = max(max_repairs, repairs)

    elapsed = time.monotonic() - start
    result = {
        "limit": LIMIT,
        "checked_nonpowers": checked,
        "precision_decimal_digits": PREC,
        "atanh_terms": TERMS,
        "max_initial_guess_repairs": max_repairs,
        "minimum_absolute_interval_endpoint": str(min_abs_endpoint),
        "elapsed_seconds": elapsed,
        "conclusion": (
            f"For every integer 2 <= A <= {LIMIT}, "
            "A^(log(3)/log(2)) is integral only if A is a power of two."
        ),
    }
    print(json.dumps(result, indent=2, sort_keys=True))


if __name__ == '__main__':
    main()
\end{lstlisting}

\begin{thebibliography}{99}

\bibitem{AlaogluErdos1944}
L.~Alaoglu and P.~Erd\H{o}s,
\newblock On highly composite and similar numbers,
\newblock \emph{Transactions of the American Mathematical Society} 56 (1944),
448--469.

\bibitem{BakerBook}
A.~Baker,
\newblock \emph{Transcendental Number Theory},
\newblock Cambridge University Press, 1975.

\bibitem{BombieriGubler}
E.~Bombieri and W.~Gubler,
\newblock \emph{Heights in Diophantine Geometry},
\newblock Cambridge University Press, 2006.

\bibitem{HardyLittlewoodPolya}
G.~H.~Hardy, J.~E.~Littlewood, and G.~P\'olya,
\newblock \emph{Inequalities}, second edition,
\newblock Cambridge University Press, 1952.

\bibitem{HarishChandra}
Harish-Chandra,
\newblock Differential operators on a semisimple Lie algebra,
\newblock \emph{American Journal of Mathematics} 79 (1957), 87--120.

\bibitem{ItzyksonZuber}
C.~Itzykson and J.-B.~Zuber,
\newblock The planar approximation. II,
\newblock \emph{Journal of Mathematical Physics} 21 (1980), 411--421.

\bibitem{Karlin}
S.~Karlin,
\newblock \emph{Total Positivity}, Vol.~I,
\newblock Stanford University Press, 1968.

\bibitem{KuipersNiederreiter}
L.~Kuipers and H.~Niederreiter,
\newblock \emph{Uniform Distribution of Sequences},
\newblock Wiley, 1974.

\bibitem{LaurentMignotteNesterenko}
M.~Laurent, M.~Mignotte, and Yu.~Nesterenko,
\newblock Formes lin\'eaires en deux logarithmes et d\'eterminants
 d'interpolation,
\newblock \emph{Journal of Number Theory} 55 (1995), 285--321.

\bibitem{LangTranscendental}
S.~Lang,
\newblock \emph{Introduction to Transcendental Numbers},
\newblock Addison--Wesley, 1966.

\bibitem{Matveev}
E.~M.~Matveev,
\newblock An explicit lower bound for a homogeneous rational linear form in
 logarithms of algebraic numbers. II,
\newblock \emph{Izvestiya: Mathematics} 64 (2000), 1217--1269.

\bibitem{Waldschmidt}
M.~Waldschmidt,
\newblock \emph{Diophantine Approximation on Linear Algebraic Groups},
\newblock Springer, 2000.

\end{thebibliography}

\end{document}
